\documentclass[11pt,a4paper]{article}
\usepackage{amsmath,amssymb,amsthm,mathtools,booktabs,microtype}
\usepackage[margin=25mm]{geometry}
\usepackage[colorlinks=true,linkcolor=blue,citecolor=blue,urlcolor=blue]{hyperref}
\setlength{\emergencystretch}{2em}
\newtheorem{theorem}{Theorem}[section]
\newtheorem{lemma}[theorem]{Lemma}
\newtheorem{proposition}[theorem]{Proposition}
\newtheorem{corollary}[theorem]{Corollary}
\newcommand{\Q}{\mathbb Q}\newcommand{\Z}{\mathbb Z}
\DeclareMathOperator{\Nm}{Nm}
\DeclareMathOperator{\rank}{rank}
\title{A ten-coordinate definition of the rational subfield\\
by a guarded finite algebra}
\author{\(\Z\Q\) Diophantine Project}
\date{19 November 2026}
\begin{document}\maketitle
\begin{abstract}
The amplified rationality detector of the preceding v81 paper can be
compressed from sixteen to ten total coordinates over the ambient
number field. Four point ordinates and two intrinsic-pole roots are
encoded in a finite algebra of rank \(64\). The proof supplies rational
point preimages, verifies every guard at all geometric characters,
uses the actual class-seven intrinsic-pole witnesses, and gives a
finite separating menu with \(10081\) entries. A shifted characteristic
determinant combines the primitive element and the inverse of every
guard in one coordinate; it does not require the primitive element
to be nonzero. The resulting integer polynomial is specified by finite
products and determinants, has content one, and has explicit degree
and coefficient bounds in the fixed circuit data. Its arithmetic
assumptions remain the six-point height inequality, positive elliptic
rank stable from \(\Q\), and Burungale's published simultaneous-twist
theorem. The exceptional-set correction is ineffective. This is a
definition of \(\Q\) inside a number field, not a reduction of the
separate variable counts for defining \(\Z\) over \(\Q\).
\end{abstract}

\section{Arithmetic assumptions and the uncompressed source}

Fix a number field \(K\). Assume that an elliptic curve \(E/\Q\) has
\[
 0<\rank E(\Q)=\rank E(K).
\]
Choose \(m\ge1\) with \(mE(K)\subset E(\Q)\); such an \(m\) exists
because the quotient of these finitely generated groups is finite.
Use an integral short Weierstrass equation
\begin{equation}\label{eq:curve}
 E:\quad y^2=F(x),\qquad F(X)=X^3+aX+b,\qquad
 a,b\in\Z,\quad 4a^3+27b^2\ne0.
\end{equation}
Put \(\xi=-x/y\) and
\[
 \Gamma^\circ=\{Q\in mE(K)\setminus\{O\}:x(Q)y(Q)\ne0\},\qquad
 \mathcal A=\{\xi(U)/\xi(V):U,V\in\Gamma^\circ\}\subset\Q^\times.
\]

The preceding paper~\cite{amplified} proves the following source
statement, with its local propagation and completeness audited
independently. Set
\begin{equation}\label{eq:f}
 f(T)=\frac{T^3-2}{T^3-3},\qquad
 \mathcal M=\{\pm1,\pm2,\pm3,\pm5,\pm6,\pm7,\pm10,\pm14\},
\end{equation}
and
\begin{equation}\label{eq:pole}
 P_d(v,s)=d(d+v^2)^2s^4-d(d-v^2)^2(1-d)^2.
\end{equation}
For \(g=f(t)\), \(\lambda,\beta\in\mathcal A\), define
\begin{equation}\label{eq:sourceinputs}
 b_1=\frac{(\lambda-g)^2(1+g)^2}{g^3},\qquad
 b_2=(\beta-t)^2\left(\frac{(1+\lambda)^2}{\lambda}\right)^3.
\end{equation}
Require \(g\ne0,\infty\), and require, for some
\(\mu_1,\mu_2\in\mathcal M\), pole equations over \(K\)
\[
                    Z_i^2=P_{b_i/\mu_i}(v_i,s_i)\qquad(i=1,2).
\]
Here all coordinates are quantified in \(K\); the rationality of
\(\lambda,\beta\) follows from rank stability.

Let \(D\) be the reduced divisor of the six roots of
\((T^3-2)(T^3-3)\). Assume the one-divisor lower bound
\begin{equation}\label{eq:height}
 N^{(1)}(D,t)\ge\theta h(t)-C\quad(t\in K\setminus D),
 \qquad \theta>7/2.
\end{equation}
Ordinary \(abc\) over \(K\) implies this estimate. The source predicate
contains every rational input and only finitely many irrational inputs.
Consequently there is \(H\in\Z[T]\), with no rational root, such that
adjoining \(H(t)\ne0\) makes this predicate define exactly \(\Q\).
One can take the product of the primitive minimal polynomials of the
irrational accepted inputs. No effective specification of \(H\) is asserted.

The pole completeness used here is the class-seven consequence of
Burungale's published Theorem~B~\cite{Burungale}. Its precise density
form is proved again in Section~\ref{sec:density}; a square-class-one
extension is not assumed. The rationality-descent architecture is due
to Hector Pasten and Xavier Vidaux~\cite{PV}. The amplification and
preparation are the preceding project work~\cite{amplified}; the
finite-algebra compression below is a further improvement. The
eleven-coordinate rational-integer theorem used only in the final
composition corollary originates in notes of \textbf{Gabe Conant with
Astra 6.2}, with the project's subsequent audits and improvements.

\begin{theorem}\label{thm:main}
Under the stated stable-rank, height, and class-seven arithmetic inputs,
there is a polynomial in \(\Z[T,X_1,X_2,X_3,X_4,v_1,s_1,v_2,s_2,Z]\)
whose input projection over \(K\) is exactly \(\Q\).
Thus the definition uses at most ten total \(K\)-coordinates,
including its distinguished input.
The finite determinant specification below has integer content one.
Its degree, support, and coefficient bounds are given in
Section~\ref{sec:bounds}.
\end{theorem}

\section{Approximation with rational point preimages}

We will use the stronger completeness option in which all four
elliptic preimages lie in \(E(\Q)\). The positive-rank subgroup
\(\Gamma_0=mE(\Q)\subseteq mE(K)\) suffices.

\begin{lemma}\label{lem:approx}
Ratios \(\xi(U)/\xi(V)\), with
\(U,V\in\Gamma_0\) in the allowed affine chart, approximate any rational
target simultaneously at the real place and at any specified finite
set of rational primes. They can avoid any fixed finite set of rational
values. The two points have rational affine preimages under \([m]\).
\end{lemma}
\begin{proof}
This is the approximation proof of~\cite{amplified}, applied to
\(\Gamma_0\) instead of \(mE(K)\). For completeness, for a nonzero
target \(a_0/b_0\), integral formal multiplication gives
\[
 \xi([a_0]Q)=a_0z+z^2F_{a_0}(z),\quad
 \xi([b_0]Q)=b_0z+z^2F_{b_0}(z),\quad z=\xi(Q).
\]
If \(v_p(z)>v_p(b_0)\), the error of the ratio has valuation at least
\(v_p(z)-2v_p(b_0)\). The subgroup lying in any chosen finite list
of deep formal neighborhoods has finite index and positive rank.
Nonzero multiples of a nontorsion point in it approach \(O\) in the
compact group \(E(\mathbb R)\). The ratio then approaches \(a_0/b_0\)
really while preserving the finite-place bounds. Zero targets and
finite avoidance follow by first choosing a nearby nonzero rational
target outside the forbidden set and shrinking its neighborhoods.
The proof may be performed in \(mE(\Q)\) throughout, so both points
have rational preimages. The deep formal points have nonzero
abscissa and ordinate.
\end{proof}

For any fixed \(t\in\Q\), the values \(f(t)=0,\infty,-1\) are impossible
by the rational-root test. Choose the first ratio \(\lambda\) close
to \(g=f(t)\) really and at all needed finite primes, avoiding
\[
                       \{0,1,-1,g,-g\}.
\]
After fixing it, choose the second ratio \(\beta\) close to \(t\),
avoiding \(\{t,-t\}\). The second finite-place preparation includes
every numerator and denominator prime of the already chosen
\(\lambda(1+\lambda)\). The proof of~\cite{amplified} therefore gives,
for both actual source inputs,
\begin{equation}\label{eq:preparedb}
 \begin{gathered}
 b_i\ne0,\quad |b_i|<1,\quad b_i\text{ has square reduced denominator},\\
 v_2(b_i)\ge5,\quad v_3(b_i),v_5(b_i),v_7(b_i)\ge1.
 \end{gathered}
\end{equation}
The additional finite forbidden values do not change that proof.
For each \(i\), a choice \(\mu_i\in\mathcal M\) makes
\begin{equation}\label{eq:preparedd}
 \begin{gathered}
 d_i=b_i/\mu_i:\qquad 0<d_i<1,\quad
 d_i\text{ has square reduced denominator},\\
 d_i\in7\Q_2^{\times2},\qquad v_2(d_i)\ge4.
 \end{gathered}
\end{equation}
Indeed, choose \(\mu_i b_i>0\) and \(-\mu_i b_i\in\Q_2^{\times2}\).
The prescribed valuations prevent an odd denominator at the primes
dividing \(\mu_i\), and \(b_i/\mu_i\) has the same square class as
\(\mu_i b_i\).

\section{The exact class-seven density input}\label{sec:density}

\begin{lemma}\label{lem:density}
For every rational \(d\) satisfying~\eqref{eq:preparedd}, the projection
to the \((v,s)\)-plane of rational solutions
\(Y^2=P_d(v,s)\) is Zariski dense, including over \(\overline\Q\).
This follows from the published case
\(\beta\equiv7\pmod8,\ \eta\equiv1\pmod8\) of Burungale's Theorem~B.
\end{lemma}
\begin{proof}
Write \(d=N/Q_0^2\) in lowest terms, with \(0<N<Q_0^2\).
Let \(\beta_0,\eta\) be the positive squarefree representatives of
\(d\) and \(1-d\). The integer numerators \(N,Q_0^2-N\) are coprime,
so \(\beta_0,\eta\) are coprime. The two-adic assumptions give
\(\beta_0\equiv7\pmod8\) and \(\eta\equiv1\pmod8\).
At every odd prime dividing \(\beta_0\), the number \(1-d\) is a
unit congruent to one; thus \(\eta\) is a local square. At every
odd prime dividing \(\eta\), the number \(d\) is a unit congruent to
one; thus \(\beta_0\) is a local square.

The cited Theorem~B supplies a positive \(n\) for which
\(E_n:y^2=x^3-n^2x\) and \(E_{n\eta}\) both have positive rank and
points with nonzero abscissa in \(\beta_0\Q^{\times2}=d\Q^{\times2}\).
Write \(1-d=\eta c^2\) and scale the second curve to
\(E_{n(1-d)}\). On each curve the abscissa square-class map is a
homomorphism into \(\Q^\times/\Q^{\times2}\); its value is unchanged
upon adding a point in \(2E_l(\Q)\). Hence each prescribed class
contains infinitely many rational points. Discard the finitely many
torsion points and choose the two points independently:
\[
 R=(d a_1^2,a_1U)\in E_n(\Q),\qquad
 R'=(d a_2^2,a_2V)\in E_{n(1-d)}(\Q).
\]
Their positive abscissas exceed the positive curve parameters, so
all the following denominators are nonzero and \(0<e<1\):
\[
 e=\frac{n}{d a_1^2},\qquad p_0=\frac{U}{d a_1^2},\qquad
 s=\frac{a_2}{a_1},\qquad q_0=\frac{V}{d a_1^2}.
\]
The curve equations give
\[
 p_0^2=d(1-e^2),\qquad
 q_0^2=d\bigl(s^4-e^2(1-d)^2\bigr).
\]
Put \(v=p_0/(1+e)\). Then
\[
 e=\frac{d-v^2}{d+v^2},\qquad
 p_0=\frac{2dv}{d+v^2},\qquad Y=(d+v^2)q_0
\]
give exactly \(Y^2=P_d(v,s)\).

As \(R\) varies there are infinitely many distinct \(a_1^2\), hence
infinitely many \(e\) and \(v^2=d(1-e)/(1+e)\). For every fixed first
point, the independently varying second point gives infinitely many
\(s\). A polynomial over \(\overline\Q\) vanishing at every such pair
vanishes as a polynomial in \(s\) for infinitely many \(v\), so every
coefficient vanishes identically. This proves the asserted density.
No assertion for the square class \(\beta_0=1\) is used.
\end{proof}

In particular, any finite list of nonzero polynomials in the two actual
pole parameters can be avoided while retaining rational witnesses.
This is an avoidance statement; it does not prescribe a new rational
square at some other character.

\section{A polynomial finite algebra in nine base coordinates}

Since \(\xi\circ[m]\) is odd under \((x,y)\mapsto(x,-y)\), choose
nonzero \(u_m,v_m\in\Z[X]\) such that
\begin{equation}\label{eq:xi-map}
              \xi([m](x,y))=y\,u_m(x)/v_m(x)
\end{equation}
as a rational function on \(E\). Divide out common polynomial factors.
This function can be computed by the usual multiplication formulas.
On the chart \(F(x)u_m(x)v_m(x)\ne0\), its image under \([m]\)
lies in \(\Gamma^\circ\) whenever the point lies in \(E(K)\).
Indeed, on a short Weierstrass model, \(\xi\) vanishes at \(O\) and
at points with \(x=0\), and has a pole at every point with
\(y=0,x\ne0\). If \((0,0)\) is on the curve, \(\xi\) has a zero
there. Thus a finite nonzero value lies in the stated chart.
Conversely, a preimage of a point in \(\Gamma^\circ\) has
\(F u_m v_m\ne0\): it cannot be a two-torsion point, and the reduced
expression has neither a zero nor a pole there.

Use the polynomial ring in nine base coordinates
\[
 \mathfrak R=\Z[T,X_1,X_2,X_3,X_4,v_1,s_1,v_2,s_2].
\]
Set \(F_i=F(X_i)\), \(u_i=u_m(X_i)\), and \(v_i^*=v_m(X_i)\);
the superscript distinguishes \(v_i^*\) from the free pole parameter.
Form the free rank-sixteen algebra
\[
 \mathfrak A_0=\mathfrak R[Y_1,Y_2,Y_3,Y_4]/
                    (Y_i^2-F_i:1\le i\le4).
\]
Its specified basis is the set of squarefree monomials in the four
\(Y_i\). Define the following elements and scalar polynomials:
\[
 \Lambda=Y_1Y_2u_1v_2^*,\quad L=F_2u_2v_1^*,\qquad
 B=Y_3Y_4u_3v_4^*,\quad M=F_4u_4v_3^*.
\]
Whenever the displayed scalar factors are nonzero,
\(\lambda=\Lambda/L\) and \(\beta=B/M\) are exactly the two
ratios of~\eqref{eq:xi-map}. Also
\[
                     S=\Lambda^2=F_1F_2u_1^2(v_2^*)^2\in\mathfrak R.
\]
With \(P=T^3-2\) and \(Q=T^3-3\), put
\begin{align}
 A_1&=(\Lambda Q-PL)^2(P+Q)^2,&
 B_1&=L^2QP^3,\label{eq:a1}\\
 A_2&=(B-TM)^2(L+\Lambda)^6\Lambda,&
 B_2&=M^2L^3S^2.\label{eq:a2}
\end{align}
Then \(b_i=A_i/B_i\) are exactly~\eqref{eq:sourceinputs}.
For the second identity use \(1/\Lambda^3=\Lambda/S^2\).
Both \(B_i\) are scalar polynomials, so no algebra-valued denominator
is hidden in these formulas.

Fix one of the \(16^2\) pairs \((\mu_1,\mu_2)\), and set
\(C_i=\mu_iB_i\). Define
\begin{equation}\label{eq:radicands}
 R_i=A_iC_i^3(A_i+C_iv_i^2)^2s_i^4
       -A_iC_i(A_i-C_iv_i^2)^2(C_i-A_i)^2\in\mathfrak A_0.
\end{equation}
The identity
\[
                     R_i=C_i^6P_{A_i/C_i}(v_i,s_i)
\]
holds whenever \(C_i\ne0\). Adjoin the two monic quadratic roots:
\[
 \mathfrak A=\mathfrak A_0[W_1,W_2]/(W_1^2-R_1,\ W_2^2-R_2).
\]
This algebra is free of rank \(64\). Its specified basis consists of
the squarefree monomials in the six generators. Under a field-valued
character, the original pole root is \(Z_i=W_i/C_i^3\).

\section{All-character guards and rational completeness}

For \(a\in\mathfrak A_0\), let \(N_0(a)\in\mathfrak R\) be the determinant
of multiplication by \(a\) in its rank-sixteen basis. For the current
multiplier pair define
\begin{equation}\label{eq:guard}
 G_*=
 H(T)P Q(P+Q)\!
 \left(\prod_{i=1}^4F_i u_i v_i^*\right)
 N_0(A_1)N_0(A_2)N_0(R_1)N_0(R_2)\in\mathfrak R.
\end{equation}
The condition \(G_*\ne0\) makes \(L,M,S,B_1,B_2,C_1,C_2\) nonzero.
It also makes \(A_1,A_2,R_1,R_2\) units in the specialized point
algebra. The four nonzero \(F_i\) make that algebra finite etale in
characteristic zero, and the two unit radicands make the full algebra
finite etale of rank \(64\).

\begin{lemma}\label{lem:guards}
For every \(t\in\Q\), there are rational values of the other eight
base coordinates, a multiplier pair, and a rational character of
the resulting algebra \(\mathfrak A\), for which \(G_*\ne0\).
\end{lemma}
\begin{proof}
Choose \(\lambda,\beta\) as in Section~2 using rational preimages.
The four chosen abscissas are rational, and the point algebra
\(\mathfrak A_0\otimes\Q\) is split: its sixteen geometric characters
independently change the signs of the four nonzero rational ordinates.
They change \(\lambda\) to \(\lambda\) or \(-\lambda\), and \(\beta\)
to \(\beta\) or \(-\beta\). Every point-chart factor in~\eqref{eq:guard}
is nonzero. Also \(H(t)P(t)Q(t)(P(t)+Q(t))\ne0\), since none of these
polynomials has a rational root.

The avoidance \(\lambda\notin\{\pm g,\pm1\}\) and
\(\beta\notin\{\pm t\}\) makes both source inputs nonzero at every
point character. Thus \(N_0(A_1)N_0(A_2)\ne0\). Choose
\(\mu_1,\mu_2\) satisfying~\eqref{eq:preparedd} at the selected
rational character.

For either pole block, every other point character sends \(A_i/C_i\)
to a nonzero rational number \(d'\). Its radicand is
\(C_i^6P_{d'}(v_i,s_i)\), a nonzero polynomial in the two free
parameters: the coefficient of \(s_i^4\) is
\(C_i^6d'(d'+v_i^2)^2\), which is not the zero polynomial.
Consequently the product over all sixteen point characters, namely
\(N_0(R_i)\), is a nonzero polynomial in these two parameters.
Apply Lemma~\ref{lem:density} to the actual prepared \(d_i\), avoiding
that product. The selected pole root remains rational and all other
radicands become nonzero. Perform this independently for the two
blocks. Every factor of \(G_*\) is now nonzero, and the selected
point and pole values give a rational character of \(\mathfrak A\).
\end{proof}

This proof requires only avoidance of nonzero polynomials at the other
characters. Those other characters need not themselves have rational
pole witnesses. For arbitrary base values over \(K\), the point
characters need not be rational either; their role in the soundness
argument is through finite-algebra units.

\section{A finite primitive menu and one guarded encoder}

For \(1\le j\le10081\), set
\[
 \alpha_j=Y_1+jY_2+j^2Y_3+j^3Y_4+j^4W_1+j^5W_2.
\]
Let \(M_j\) be its \(64\times64\) multiplication matrix in the specified
basis of \(\mathfrak A\). Let
\[
 \Delta_j=\det[\,1,\alpha_j,\alpha_j^2,\ldots,\alpha_j^{63}\,]
                 \in\mathfrak R,
\]
where the bracket denotes the matrix of coefficient columns in that
basis. Define \(G_j=G_*\Delta_j\) and the single-coordinate polynomial
\begin{equation}\label{eq:encoder}
           E_{\mu_1,\mu_2,j}(Z)
           =\det\bigl(G_j(ZI_{64}-M_j)-I_{64}\bigr)
           \in\mathfrak R[Z].
\end{equation}

\begin{lemma}\label{lem:separator}
At every characteristic-zero specialization with \(G_*\ne0\), at least
one \(j\in\{1,\ldots,10081\}\) has \(\Delta_j\ne0\).
\end{lemma}
\begin{proof}
The specialized algebra is etale of rank \(64\), so it has \(64\)
distinct geometric characters. For any two different characters,
their values differ on at least one of its six named generators.
The difference of their values on
\(Y_1+JY_2+\cdots+J^5W_2\) is therefore a nonzero polynomial in \(J\)
of degree at most five. It vanishes at at most five integers.
There are \(\binom{64}{2}=2016\) pairs of characters. Hence at most
\(5\cdot2016=10080\) integers fail to separate some pair.
The menu of \(10081\) distinct integers contains a separator.

For such a separator, the evaluation matrix of
\(1,\alpha_j,\ldots,\alpha_j^{63}\) is a Vandermonde matrix with
nonzero determinant. The evaluation matrix of any basis of an etale
algebra is invertible after passing to a splitting field. Therefore
\(\Delta_j\ne0\).
\end{proof}

\begin{lemma}\label{lem:encoder}
For a specialization of the nine base coordinates in \(K\), a root
in \(K\) of~\eqref{eq:encoder} exists if and only if
\(G_j\ne0\) and the specialized algebra has a \(K\)-valued character.
\end{lemma}
\begin{proof}
If \(G_j=0\), the determinant is \(\det(-I_{64})=1\).
Thus every root forces \(G_j\ne0\), and hence \(\Delta_j\ne0\).
The powers \(1,\alpha_j,\ldots,\alpha_j^{63}\) are then a basis.
The minimal polynomial of \(\alpha_j\) has degree \(64\) and equals
its multiplication characteristic polynomial \(p_j\). The algebra is
isomorphic to \(K[A]/(p_j(A))\).

When \(G_j\ne0\), equation~\eqref{eq:encoder} is
\[
              G_j^{64}p_j(Z-1/G_j)=0.
\]
Its root supplies the \(K\)-valued character
\(\alpha_j\mapsto Z-1/G_j\). Conversely, if \(\chi\) is a
\(K\)-valued character, take \(Z=\chi(\alpha_j)+1/G_j\).
No nonzero condition on \(\chi(\alpha_j)\) is required.
\end{proof}

The shift by the guard inverse is essential to this formula. Replacing
it without proof by a determinant of \(ZG_jM_j-I\) would impose an
additional nonvanishing condition on the selected primitive value.

\section{The complete polynomial and source equivalence}

Define
\begin{equation}\label{eq:final}
 \mathcal F(T,X_1,X_2,X_3,X_4,v_1,s_1,v_2,s_2,Z)
  =\prod_{(\mu_1,\mu_2)\in\mathcal M^2}
       \ \prod_{j=1}^{10081}E_{\mu_1,\mu_2,j}(Z).
\end{equation}
All factors reuse the same ten coordinates. The algebras, guards,
and multiplication matrices in a factor are constructed using that
factor's multiplier pair.

\begin{proof}[Proof of Theorem~\ref{thm:main}: soundness]
Suppose \(\mathcal F=0\) at a tuple in \(K^{10}\). One factor vanishes.
Lemma~\ref{lem:encoder} gives a \(K\)-valued character \(\chi\) and
all guards in \(G_*\). The four points
\[
                       (X_i,\chi(Y_i))\in E(K)
\]
lie in the chart of~\eqref{eq:xi-map}. Their images under \([m]\)
belong to \(\Gamma^\circ\subset E(\Q)\), so the recovered
\(\lambda=\chi(\Lambda)/L\) and \(\beta=\chi(B)/M\) are in
\(\mathcal A\subset\Q^\times\).

Equations~\eqref{eq:a1}--\eqref{eq:a2} recover the exact source inputs
\(b_i=\chi(A_i)/B_i\). Since \(C_i\ne0\), the roots
\(\chi(W_i)/C_i^3\) satisfy precisely
\(Z_i^2=P_{b_i/\mu_i}(v_i,s_i)\).
The contact denominators and \(H(T)\) are nonzero. Thus this tuple
gives an accepted tuple of the corrected source predicate, whose
input belongs to \(\Q\).
\end{proof}

\begin{proof}[Proof of Theorem~\ref{thm:main}: completeness]
For \(t\in\Q\), Lemma~\ref{lem:guards} gives rational values of the
other eight base coordinates and a multiplier pair for which \(G_*\ne0\)
and the specialized algebra has a rational character \(\chi\).
Lemma~\ref{lem:separator} supplies a \(j\) with \(\Delta_j\ne0\).
Then \(Z=\chi(\alpha_j)+1/G_j\) is rational, and its factor in
\eqref{eq:final} vanishes by Lemma~\ref{lem:encoder}.
This proves completeness with all witnesses in \(\Q\subset K\).
\end{proof}

No inverse coordinate besides the encoder \(Z\) has been used.
The six monic algebraic generators are not free coordinates. The
complete census is
\[
 \underbrace{1}_{\text{input}}+
 \underbrace{4}_{\text{point abscissas}}+
 \underbrace{4}_{\text{pole parameters}}+
 \underbrace{1}_{\text{encoder}}=10.
\]

\section{Degree, support, and coefficient normalization}\label{sec:bounds}

The finite products and matrices above are a literal specification;
they need not be expanded to define the polynomial. Put
\[
 n_{\rm fac}=16^2\cdot10081=2\,580\,736.
\]
Let \(d_M\) be a common upper bound for the total degrees, in the nine
base coordinates, of all entries in the six generator multiplication
matrices, over all multiplier pairs. Let \(h_*\) bound the degrees
of all guards \(G_*\). These are finite circuit data determined by
the chosen multiplication map, \(H\), and the displayed formulas.

The entries in \(M_j\) have degree at most \(d_M\).
The coefficient column of \(\alpha_j^r\) has degree at most \(rd_M\);
therefore
\[
 \deg\Delta_j\le\left(\sum_{r=0}^{63}r\right)d_M=2016d_M.
\]
Consequently a valid bound for the total degree of~\eqref{eq:final} is
\begin{equation}\label{eq:degree}
 D_{\rm tot}
 =n_{\rm fac}\,64\bigl(h_*+2016d_M+\max\{1,d_M\}\bigr).
\end{equation}
Its support has cardinality at most \(\binom{D_{\rm tot}+10}{10}\).
These are bounds, not claims of exact degree or support.

Every multiplication matrix and every guard has integer polynomial
entries. Each factor in~\eqref{eq:encoder} is primitive. Indeed, modulo
any rational prime, if \(G_j\) becomes the zero polynomial the factor
is the constant one; otherwise its leading coefficient in \(Z\) is
the nonzero polynomial \(G_j^{64}\). Thus no prime divides every
coefficient of a factor. Gauss's lemma makes the product primitive.

For an explicit coefficient bound in the same data, write
\(\|A\|_1\) for the sum of absolute values of the coefficients of an
integer polynomial. Let
\[
 B_*=\max_{\mu_1,\mu_2,j}
 \left(1+\|G_j\|_1
       \left(1+\max_{r,s}\|(M_j)_{rs}\|_1\right)\right).
\]
The determinant expansion gives
\[
 \|\mathcal F\|_1
 \le \left(64!\,B_*^{64}\right)^{n_{\rm fac}}.
\]
In particular the coefficient height is at most this bound.
No effective numerical values for the exceptional correction \(H\),
the rank-stability multiplier, or these circuit bounds are claimed.

\section{Composition and limits}

\begin{corollary}
If an \(N\)-coordinate rational polynomial defines \(\Z\) over \(\Q\),
then, under the hypotheses of Theorem~\ref{thm:main}, \(\Z\) has a
definition over \(K\) in \(10N\) total \(K\)-coordinates.
Its ring of integers has a definition in
\(10N[K:\Q]\) total \(K\)-coordinates.
The project's eleven-coordinate theorem therefore gives \(110\)
and \(110[K:\Q]\), with all of that theorem's retained premises.
\end{corollary}
\begin{proof}
Restrict each of the \(N\) original coordinates, including the input,
by a separate copy of the ten-coordinate selector. This uses
\(N+9N=10N\) coordinates. For the ring of integers, expand in an
integral basis and eliminate the last coefficient using the
distinguished input, as in~\cite{amplified}. This multiplies the count
by \([K:\Q]\).
\end{proof}

For quadratic \(K\), Zywina's Theorem~1.2~\cite{Zywina}, with base
field \(\Q\), supplies the stable rank-one elliptic curve. Thus the
displayed counts become \(110\) for \(\Z\) in \(K\) and \(220\) for
\(\mathcal O_K\), under the stated height and arithmetic inputs.
They are field-valued existence bounds, not rational-variable records,
quartic bounds, or effective coefficient lists.

The compression changes how six finite roots and all guards are encoded.
It does not replace either free intrinsic-pole parameter by a prescribed
rational equality. In particular it does not solve the separate
realization problem needed to lower the \(25/11\) counts over \(\Q\).

There is also a precise limit on the contact-map degree within the
stated amplification criterion. A rational map of degree at most two
has at most four distinct odd-order zeros and poles, including infinity.
Thus \(r\le4\), and the ordinary-\(abc\) coefficient
\(r-4-2\epsilon\) cannot exceed \(2/k\) for any finite \(k\ge1\).
The cubic map is therefore minimal for this particular criterion.
This is not a lower bound for all possible rationality constructions.

\section{Independent verification}

The independent compression checker reconstructs the source input
identities and the pole-root clearing, verifies the separating bound
and coordinate census, and checks the encoder on exact finite-algebra
fixtures. The accompanying review separately checks rational preimages,
the published class-seven density input, all-character guards, and the
two directions of source equivalence. A finite fixture does not prove
the universal primitive-element argument or the arithmetic premises;
those proofs are given above and in the cited source.

\begin{thebibliography}{9}\footnotesize\raggedright
\setlength{\itemsep}{1pt}
\bibitem{amplified} \(\Z\Q\) Diophantine Project,
\emph{Amplified rationality descent: a cubic rational map and a
sixteen-coordinate selector}, 19 November 2026,
\nolinkurl{papers/number_fields/amplified_rationality_descent_20261119_v81.tex}.
\bibitem{PV} H.~Pasten and X.~Vidaux,
\emph{On Hilbert's tenth problem for number fields},
arXiv:2611.24748v1, 18 November 2026,
\url{https://arxiv.org/html/2611.24748v1}.
\bibitem{Burungale} A.~Burungale,
\emph{A Diophantine definition of \(\Z\) in \(\Q\)},
arXiv:2611.16130v1, Theorem~B,
\url{https://arxiv.org/html/2611.16130v1}.
\bibitem{intrinsic} \(\Z\Q\) Diophantine Project,
\emph{A two-coordinate pole test of coefficient height four},
17 November 2026,
\nolinkurl{papers/pole_parameters/intrinsic_pole_20261117_v56.tex}.
\bibitem{Zywina} D.~Zywina,
\emph{Rank one elliptic curves and rank stability},
arXiv:2505.16960v1, Theorem~1.2,
\url{https://arxiv.org/html/2505.16960v1}.
\end{thebibliography}
\end{document}
